\chapter{Soft Differentiable Resampling}
\label{ch:bf-dpf-soft-resampling}

The previous chapters established the classical particle-filter setting in which a
weighted cloud of particles is turned back into an equally weighted cloud by a
resampling step.  The mathematical difficulty is that the standard resampling map
is discrete: it duplicates some particles, deletes others, and changes abruptly
when weights cross threshold values.  That makes naive pathwise differentiation
break down.

This chapter studies the simplest smooth replacement of that discrete step.  It
does not introduce optimal transport yet.  Its narrower task is to show, in the
smallest possible setting, how differentiability can be bought by changing the
resampling map itself.

\section{What problem this chapter solves}
\label{sec:bf-soft-what-problem}

Suppose a weighted particle cloud is
\begin{equation}
\label{eq:bf-soft-weighted-cloud}
  \hat\pi_t^N(dx)
  =
  \sum_{i=1}^N w_t^{(i)}\,\delta_{x_t^{(i)}}(dx),
  \qquad
  w_t^{(i)}\ge 0,
  \quad
  \sum_{i=1}^N w_t^{(i)}=1.
\end{equation}
Classical multinomial resampling tries to convert this into an equally weighted
cloud by drawing ancestor indices from the categorical law defined by the weights.
This is fine as a simulation rule.  It is not fine if one wants the realized
output cloud to depend smoothly on the weights.

The chapter therefore asks one very specific question:
\begin{quote}
How can we make the resampled cloud depend more smoothly on the weights, and what
bias is introduced when we do so?
\end{quote}

\section{Running two-particle discontinuity cell}
\label{sec:bf-soft-two-particle-cell}

The smallest useful example has two particles.  Let
\[
  x_1=-1,
  \qquad
  x_2=2,
  \qquad
  w=(p,1-p),
  \qquad
  p\in(0,1),
\]
and let one fixed common random number \(U\in(0,1)\) be used to define the first
resampled ancestor.  Standard categorical resampling chooses
\begin{equation}
\label{eq:bf-soft-categorical-choice}
  a(U,p)
  =
  \begin{cases}
    1, & U\le p,\\
    2, & U>p.
  \end{cases}
\end{equation}
For a fixed \(U\), this is a step function of \(p\): small changes in the weight
usually do nothing, and then a threshold crossing changes the realized ancestor
abruptly.

\paragraph{Plain-language takeaway.}
If the algorithm literally copies an ancestor chosen by a discrete threshold, the
output particle usually stays exactly the same when the weights move a little.  A
pathwise derivative of that realized output is therefore typically zero until a
jump occurs.

\section{Classical resampling bottleneck}
\label{sec:bf-soft-bottleneck}

For standard multinomial resampling, the equal-weight cloud is obtained by
sampling ancestor indices
\[
  a_j\sim\mathrm{Categorical}(w_t^{(1:N)})
\]
and setting
\begin{equation}
\label{eq:bf-soft-hard-resample}
  \tilde x_t^{(j)} = x_t^{(a_j)}.
\end{equation}
This map is exact for the classical resampling law, but it is not pathwise
differentiable with respect to the weights because the realized ancestor path is
piecewise constant in those weights.

\section{Soft resampling rule}
\label{sec:bf-soft-rule}

A simple smooth surrogate interpolates between the sampled ancestor and the
weighted mean of the cloud.  Define
\begin{equation}
\label{eq:bf-soft-mean}
  \bar x_t = \sum_{i=1}^N w_t^{(i)}x_t^{(i)},
\end{equation}
and then set
\begin{equation}
\label{eq:bf-soft-rule-eqn}
  \tilde x_t^{(j)}
  =
  \alpha x_t^{(a_j)} + (1-\alpha)\bar x_t,
  \qquad
  \alpha\in[0,1].
\end{equation}
When \(\alpha=1\), the rule reduces to hard categorical resampling.  When
\(\alpha<1\), the output depends smoothly on the weighted mean, so the realized
output cloud becomes less discontinuous as a function of the weights.

\section{Worked two-particle calculation}
\label{sec:bf-soft-worked-cell}

In the two-particle cell,
\[
  \bar x = p(-1)+(1-p)2 = 2-3p.
\]
The soft-resampling output for one realized ancestor is therefore
\begin{align}
\label{eq:bf-soft-two-particle-output}
  \tilde x(U,p)
  &=
  \alpha x_{a(U,p)} + (1-\alpha)(2-3p).
\end{align}
Two immediate facts are now visible.

First, when the ancestor is held fixed, the dependence on \(p\) is smooth through
the weighted mean term.  Second, the object being computed has changed: the
surrogate particle is no longer equal to the chosen ancestor itself unless
\(\alpha=1\).

The weighted mean is preserved in conditional expectation because
\begin{equation}
\label{eq:bf-soft-mean-preservation}
  \mathbb E[\tilde x_t^{(j)}\mid X,w]
  =
  \alpha\sum_i w_i x_i + (1-\alpha)\bar x_t
  =
  \bar x_t.
\end{equation}
So the surrogate keeps the first-moment center of mass, but not necessarily other
nonlinear summaries of the cloud.

\section{Algorithm: soft differentiable resampling}
\label{sec:bf-soft-algorithm}

\begin{algorithm}[ht]
\caption{\texttt{soft\_resampling\_step}}
\begin{algorithmic}[1]
\Require weighted particle cloud \(\{x_i,w_i\}_{i=1}^N\), interpolation parameter \(\alpha\in[0,1]\), realized ancestor indices \(a_j\) if hard ancestors are retained, or a separately declared softened ancestor rule if they are not
\Ensure equal-weight surrogate cloud \(\widetilde X=(\tilde x^{(1)},\ldots,\tilde x^{(N)})\) and the declared ancestor policy
\State Compute the weighted mean \(\bar x = \sum_{i=1}^N w_i x_i\)
\For{each output index \(j=1,\ldots,N\)}
  \State Form \(\tilde x^{(j)} = \alpha x_{a_j} + (1-\alpha)\bar x\)
\EndFor
\State Return the equal-weight cloud \(\widetilde X\), the chosen \(\alpha\), and the declared ancestor policy used to define \(a_j\)
\end{algorithmic}
\end{algorithm}

\paragraph{Approximation enters here.}
This algorithm is not another notation for categorical resampling.  It is a new
surrogate object that trades exactness for smoother dependence on the weights.

\section{Filtering-loop insertion point}
\label{sec:bf-soft-loop}

Inside the filtering loop, the soft surrogate sits exactly where hard resampling
would have sat:
\begin{equation}
\label{eq:bf-soft-loop}
  \{x_i,w_i\}_{i=1}^N
  \xrightarrow{\mathrm{soft\ resampling}}
  \{\tilde x^{(j)},1/N\}_{j=1}^N.
\end{equation}
It changes only the equal-weighting rule.  It does not, by itself, change the
preceding weighting rule or add an OT correction.

\section{Verification contract}
\label{sec:bf-soft-verification}

A minimal verification pass should check:
\begin{enumerate}[label=(\arabic*)]
  \item the weighted mean is preserved as in
  \eqref{eq:bf-soft-mean-preservation},
  \item the surrogate becomes the hard categorical rule when \(\alpha=1\),
  \item nonlinear test-function bias is visible when expected,
  \item the declared ancestor policy is recorded explicitly,
  \item any downstream gradient claim states whether ancestors remain discrete or
  are themselves relaxed.
\end{enumerate}

A failed mean-preservation check is evidence that the surrogate has been
implemented incorrectly.  A successful mean-preservation check is not evidence
that the classical categorical target has been preserved.

\section{What this chapter does and does not claim}
\label{sec:bf-soft-boundary}

This chapter teaches the simplest smooth resampling surrogate and its bias.  It
does \emph{not} claim:
\begin{itemize}
  \item that the surrogate preserves the exact categorical resampling law,
  \item that preserving the weighted mean preserves the full distribution,
  \item that the resulting scalar is automatically the same target later used by
  HMC or training.
\end{itemize}
Its role is narrower: it shows, in the smallest possible setting, how
resampling can be made smoother and where the resulting approximation enters.

\FloatBarrier
